\section{Introduction}\label{sec:introduction}

Constructive control of a non-Abelian lattice gauge theory requires several
operations to coexist: gauge redundancy must be removed without changing the
measure, weak curvature must be represented in a chart with quantified
nonlinearity, retained variables must be separated from variables to be
integrated out, and the integration must preserve a useful effective-action
bound.  Each operation is elementary in a sufficiently small finite model.
Their simultaneous compatibility is not.  A numerical eigenvalue or a sampled
derivative maximum cannot carry a rigorous renormalization argument, while a
purely qualitative local chart does not provide a reusable quantitative input.

The heat-kernel plaquette action is a natural setting for this question.  It is
an exact compact-group interaction with a Gaussian small-time core
\cite{MenottiOnofri1981}; it is also closely related to Villain-type lattice
interactions \cite{ChevyrevGarban2025}.  The technical question answered here
is deliberately finite:

\begin{quote}
Can one exhibit a four-dimensional non-Abelian block on which gauge reduction,
retained--fiber separation, nonlinear heat-kernel control, and one exact
conditional marginalization are all joined by a machine-checkable theorem?
\end{quote}

We answer this question for the side-two four-cube and $\SUtwo$.  The gauge
choice is a rooted maximal tree.  This matters: Haar invariance makes the
quotient factorization exact and avoids an unaccounted Faddeev--Popov density.
The retained variable consists of seventeen independent linear combinations
of the twenty-four oriented boundary-face curvature sums.  It is a local
coordinate on the reduced logarithmic field, not a claim that a complete
renormalization-group block map has been constructed.

The central result, Theorem~\ref{thm:main}, has four parts.  First, exact
cubical arithmetic gives dimensions $136=17+119$, or $408=51+357$ after the
color tensor product.  Second, a rational $LDL^{\mathsf T}$ factorization gives
the cotree quadratic form the spectral floor $1/25$.  Third, deterministic
product--log and heat-kernel bounds give a scaled nonlinear Hessian remainder
below $1/400$.  Fourth, Brascamp--Lieb covariance control converts these
pointwise bounds into an exact statement about the normalized conditional
fiber law and its effective retained action.

The numerical parameters are intentionally conservative.  The point is not
that the radius $10^{-5}$ or time $10^{-18}$ is physically distinguished.  The
point is that every strict inequality has an exact rational witness and that
the whole implication can be rerun independently.  This separates a finite
mathematical fact from the much harder problem of finding a large-field
decomposition and a scale-uniform domain on which such facts can iterate.

\paragraph{Contribution and scope.}
Tree gauge, cubical Hodge theory, heat-kernel image formulas, strong
log-concavity, and Schur-complement estimates are established tools.  We do not
claim any one of them as new.  The contribution is their integration into one
explicit quotient--fiber certificate with exact ranks, a rational spectral
proof, analytic nonlinear bounds, and two independent arithmetic kernels.
This is a claim of explicit integration and certification, not priority over
constructive renormalization-group theory.

The theorem is finite-volume and local-cylinder.  Its reference Gaussian
complement estimate is not a complement bound for the interacting gauge law.
No conclusion is drawn about a volume-uniform basin, repeated blocking,
Osterwalder--Schrader reconstruction, the continuum limit, or the Hamiltonian
mass gap in the Clay problem \cite{JaffeWitten2000}.

Section~\ref{sec:block} defines the reduced block and retained map.
Section~\ref{sec:linear} proves the exact linear certificate.
Section~\ref{sec:nonlinear} controls the nonlinear action.
Section~\ref{sec:marginal} proves the conditional and effective-action
statements.  Section~\ref{sec:certificate} describes the proof artifact, and
the appendices contain all deferred estimates and reproducibility details.
